\begin{figure}[!htbp]
\centering
\begingroup
\newcommand{\nineRoot}[4]{%
\begin{tikzpicture}[x=1mm,y=1mm,baseline=0mm]
  \path[use as bounding box] (-12,-5.4) rectangle (12,4.3);
  \draw[marked edge,line width=1pt] (-2,1.4)--(2,1.4);
  \foreach \j in {1,...,#1}{
    \pgfmathsetmacro{\fy}{4.8*(\j-(#1+1)/2)/(#1-1)+1.4}
    \draw[black!65,line width=.35pt] (-2,1.4)--(-8,\fy);
    \fill (-8,\fy) circle[radius=.25mm];}
  \foreach \j in {1,...,#3}{
    \pgfmathsetmacro{\fy}{4.8*(\j-(#3+1)/2)/(#3-1)+1.4}
    \draw[black!65,line width=.35pt] (2,1.4)--(8,\fy);
    \fill (8,\fy) circle[radius=.25mm];}
  \ifnum#2>0
    \foreach \j in {1,...,#2}{
      \pgfmathsetmacro{\fx}{1.8*(\j-(#2+1)/2)}
      \draw[black!65,line width=.35pt] (-2,1.4)--(\fx,-1.9)--(2,1.4);
      \fill (\fx,-1.9) circle[radius=.25mm];}
  \fi
  \fill (-2,1.4) circle[radius=.38mm] (2,1.4) circle[radius=.38mm];
  \node[inner sep=0pt,font=\scriptsize] at (0,-4.1) {$(#1,#2,#3)^{#4}$};
\end{tikzpicture}}
\begin{tabular}{@{}r@{\hspace{1mm}}c@{\hspace{1mm}}c@{\hspace{1mm}}c@{}}
$q=11$ & \nineRoot{4}{0}{7}{2} & \nineRoot{4}{1}{6}{2} & \nineRoot{4}{2}{5}{2}\\
       & \nineRoot{4}{3}{4}{2} & \nineRoot{5}{0}{6}{2,3} & \nineRoot{5}{1}{5}{2,3}\\[.6mm]
$q=12$ & \nineRoot{5}{0}{7}{2} & \nineRoot{5}{1}{6}{2} & \nineRoot{5}{2}{5}{2}\\
       & \nineRoot{6}{0}{6}{2,3} & & \\[.6mm]
$q=13$ & \nineRoot{6}{0}{7}{2} & \nineRoot{6}{1}{6}{2} &\\[.6mm]
$q=14$ & \nineRoot{7}{0}{7}{2} & &
\end{tabular}
\endgroup
\caption{Neighborhoods of the endpoints of a repeated marked edge.
The superscript gives the multiplicities allowed by the degree bound.}
\label{fig:nine-root-profiles}
\end{figure}
